\documentclass[11pt]{article}

    \usepackage[breakable]{tcolorbox}
    \usepackage{parskip} % Stop auto-indenting (to mimic markdown behaviour)
    

    % Basic figure setup, for now with no caption control since it's done
    % automatically by Pandoc (which extracts ![](path) syntax from Markdown).
    \usepackage{graphicx}
    % Keep aspect ratio if custom image width or height is specified
    \setkeys{Gin}{keepaspectratio}
    % Maintain compatibility with old templates. Remove in nbconvert 6.0
    \let\Oldincludegraphics\includegraphics
    % Ensure that by default, figures have no caption (until we provide a
    % proper Figure object with a Caption API and a way to capture that
    % in the conversion process - todo).
    \usepackage{caption}
    \DeclareCaptionFormat{nocaption}{}
    \captionsetup{format=nocaption,aboveskip=0pt,belowskip=0pt}

    \usepackage{float}
    \floatplacement{figure}{H} % forces figures to be placed at the correct location
    \usepackage{xcolor} % Allow colors to be defined
    \usepackage{enumerate} % Needed for markdown enumerations to work
    \usepackage{geometry} % Used to adjust the document margins
    \usepackage{amsmath} % Equations
    \usepackage{amssymb} % Equations
    \usepackage{textcomp} % defines textquotesingle
    % Hack from http://tex.stackexchange.com/a/47451/13684:
    \AtBeginDocument{%
        \def\PYZsq{\textquotesingle}% Upright quotes in Pygmentized code
    }
    \usepackage{upquote} % Upright quotes for verbatim code
    \usepackage{eurosym} % defines \euro

    \usepackage{iftex}
    \ifPDFTeX
        \usepackage[T1]{fontenc}
        \IfFileExists{alphabeta.sty}{
              \usepackage{alphabeta}
          }{
              \usepackage[mathletters]{ucs}
              \usepackage[utf8x]{inputenc}
          }
    \else
        \usepackage{fontspec}
        \usepackage{unicode-math}
    \fi

    \usepackage{fancyvrb} % verbatim replacement that allows latex
    \usepackage{grffile} % extends the file name processing of package graphics
                         % to support a larger range
    \makeatletter % fix for old versions of grffile with XeLaTeX
    \@ifpackagelater{grffile}{2019/11/01}
    {
      % Do nothing on new versions
    }
    {
      \def\Gread@@xetex#1{%
        \IfFileExists{"\Gin@base".bb}%
        {\Gread@eps{\Gin@base.bb}}%
        {\Gread@@xetex@aux#1}%
      }
    }
    \makeatother
    \usepackage[Export]{adjustbox} % Used to constrain images to a maximum size
    \adjustboxset{max size={0.9\linewidth}{0.9\paperheight}}

    % The hyperref package gives us a pdf with properly built
    % internal navigation ('pdf bookmarks' for the table of contents,
    % internal cross-reference links, web links for URLs, etc.)
    \usepackage{hyperref}
    % The default LaTeX title has an obnoxious amount of whitespace. By default,
    % titling removes some of it. It also provides customization options.
    \usepackage{titling}
    \usepackage{longtable} % longtable support required by pandoc >1.10
    \usepackage{booktabs}  % table support for pandoc > 1.12.2
    \usepackage{array}     % table support for pandoc >= 2.11.3
    \usepackage{calc}      % table minipage width calculation for pandoc >= 2.11.1
    \usepackage[inline]{enumitem} % IRkernel/repr support (it uses the enumerate* environment)
    \usepackage[normalem]{ulem} % ulem is needed to support strikethroughs (\sout)
                                % normalem makes italics be italics, not underlines
    \usepackage{soul}      % strikethrough (\st) support for pandoc >= 3.0.0
    \usepackage{mathrsfs}
    

    
    % Colors for the hyperref package
    \definecolor{urlcolor}{rgb}{0,.145,.698}
    \definecolor{linkcolor}{rgb}{.71,0.21,0.01}
    \definecolor{citecolor}{rgb}{.12,.54,.11}

    % ANSI colors
    \definecolor{ansi-black}{HTML}{3E424D}
    \definecolor{ansi-black-intense}{HTML}{282C36}
    \definecolor{ansi-red}{HTML}{E75C58}
    \definecolor{ansi-red-intense}{HTML}{B22B31}
    \definecolor{ansi-green}{HTML}{00A250}
    \definecolor{ansi-green-intense}{HTML}{007427}
    \definecolor{ansi-yellow}{HTML}{DDB62B}
    \definecolor{ansi-yellow-intense}{HTML}{B27D12}
    \definecolor{ansi-blue}{HTML}{208FFB}
    \definecolor{ansi-blue-intense}{HTML}{0065CA}
    \definecolor{ansi-magenta}{HTML}{D160C4}
    \definecolor{ansi-magenta-intense}{HTML}{A03196}
    \definecolor{ansi-cyan}{HTML}{60C6C8}
    \definecolor{ansi-cyan-intense}{HTML}{258F8F}
    \definecolor{ansi-white}{HTML}{C5C1B4}
    \definecolor{ansi-white-intense}{HTML}{A1A6B2}
    \definecolor{ansi-default-inverse-fg}{HTML}{FFFFFF}
    \definecolor{ansi-default-inverse-bg}{HTML}{000000}

    % common color for the border for error outputs.
    \definecolor{outerrorbackground}{HTML}{FFDFDF}

    % commands and environments needed by pandoc snippets
    % extracted from the output of `pandoc -s`
    \providecommand{\tightlist}{%
      \setlength{\itemsep}{0pt}\setlength{\parskip}{0pt}}
    \DefineVerbatimEnvironment{Highlighting}{Verbatim}{commandchars=\\\{\}}
    % Add ',fontsize=\small' for more characters per line
    \newenvironment{Shaded}{}{}
    \newcommand{\KeywordTok}[1]{\textcolor[rgb]{0.00,0.44,0.13}{\textbf{{#1}}}}
    \newcommand{\DataTypeTok}[1]{\textcolor[rgb]{0.56,0.13,0.00}{{#1}}}
    \newcommand{\DecValTok}[1]{\textcolor[rgb]{0.25,0.63,0.44}{{#1}}}
    \newcommand{\BaseNTok}[1]{\textcolor[rgb]{0.25,0.63,0.44}{{#1}}}
    \newcommand{\FloatTok}[1]{\textcolor[rgb]{0.25,0.63,0.44}{{#1}}}
    \newcommand{\CharTok}[1]{\textcolor[rgb]{0.25,0.44,0.63}{{#1}}}
    \newcommand{\StringTok}[1]{\textcolor[rgb]{0.25,0.44,0.63}{{#1}}}
    \newcommand{\CommentTok}[1]{\textcolor[rgb]{0.38,0.63,0.69}{\textit{{#1}}}}
    \newcommand{\OtherTok}[1]{\textcolor[rgb]{0.00,0.44,0.13}{{#1}}}
    \newcommand{\AlertTok}[1]{\textcolor[rgb]{1.00,0.00,0.00}{\textbf{{#1}}}}
    \newcommand{\FunctionTok}[1]{\textcolor[rgb]{0.02,0.16,0.49}{{#1}}}
    \newcommand{\RegionMarkerTok}[1]{{#1}}
    \newcommand{\ErrorTok}[1]{\textcolor[rgb]{1.00,0.00,0.00}{\textbf{{#1}}}}
    \newcommand{\NormalTok}[1]{{#1}}

    % Additional commands for more recent versions of Pandoc
    \newcommand{\ConstantTok}[1]{\textcolor[rgb]{0.53,0.00,0.00}{{#1}}}
    \newcommand{\SpecialCharTok}[1]{\textcolor[rgb]{0.25,0.44,0.63}{{#1}}}
    \newcommand{\VerbatimStringTok}[1]{\textcolor[rgb]{0.25,0.44,0.63}{{#1}}}
    \newcommand{\SpecialStringTok}[1]{\textcolor[rgb]{0.73,0.40,0.53}{{#1}}}
    \newcommand{\ImportTok}[1]{{#1}}
    \newcommand{\DocumentationTok}[1]{\textcolor[rgb]{0.73,0.13,0.13}{\textit{{#1}}}}
    \newcommand{\AnnotationTok}[1]{\textcolor[rgb]{0.38,0.63,0.69}{\textbf{\textit{{#1}}}}}
    \newcommand{\CommentVarTok}[1]{\textcolor[rgb]{0.38,0.63,0.69}{\textbf{\textit{{#1}}}}}
    \newcommand{\VariableTok}[1]{\textcolor[rgb]{0.10,0.09,0.49}{{#1}}}
    \newcommand{\ControlFlowTok}[1]{\textcolor[rgb]{0.00,0.44,0.13}{\textbf{{#1}}}}
    \newcommand{\OperatorTok}[1]{\textcolor[rgb]{0.40,0.40,0.40}{{#1}}}
    \newcommand{\BuiltInTok}[1]{{#1}}
    \newcommand{\ExtensionTok}[1]{{#1}}
    \newcommand{\PreprocessorTok}[1]{\textcolor[rgb]{0.74,0.48,0.00}{{#1}}}
    \newcommand{\AttributeTok}[1]{\textcolor[rgb]{0.49,0.56,0.16}{{#1}}}
    \newcommand{\InformationTok}[1]{\textcolor[rgb]{0.38,0.63,0.69}{\textbf{\textit{{#1}}}}}
    \newcommand{\WarningTok}[1]{\textcolor[rgb]{0.38,0.63,0.69}{\textbf{\textit{{#1}}}}}
    \makeatletter
    \newsavebox\pandoc@box
    \newcommand*\pandocbounded[1]{%
      \sbox\pandoc@box{#1}%
      % scaling factors for width and height
      \Gscale@div\@tempa\textheight{\dimexpr\ht\pandoc@box+\dp\pandoc@box\relax}%
      \Gscale@div\@tempb\linewidth{\wd\pandoc@box}%
      % select the smaller of both
      \ifdim\@tempb\p@<\@tempa\p@
        \let\@tempa\@tempb
      \fi
      % scaling accordingly (\@tempa < 1)
      \ifdim\@tempa\p@<\p@
        \scalebox{\@tempa}{\usebox\pandoc@box}%
      % scaling not needed, use as it is
      \else
        \usebox{\pandoc@box}%
      \fi
    }
    \makeatother

    % Define a nice break command that doesn't care if a line doesn't already
    % exist.
    \def\br{\hspace*{\fill} \\* }
    % Math Jax compatibility definitions
    \def\gt{>}
    \def\lt{<}
    \let\Oldtex\TeX
    \let\Oldlatex\LaTeX
    \renewcommand{\TeX}{\textrm{\Oldtex}}
    \renewcommand{\LaTeX}{\textrm{\Oldlatex}}
    % Document parameters
    % Document title
    \title{Untitled}
    
    
    
    
    
    
    
% Pygments definitions
\makeatletter
\def\PY@reset{\let\PY@it=\relax \let\PY@bf=\relax%
    \let\PY@ul=\relax \let\PY@tc=\relax%
    \let\PY@bc=\relax \let\PY@ff=\relax}
\def\PY@tok#1{\csname PY@tok@#1\endcsname}
\def\PY@toks#1+{\ifx\relax#1\empty\else%
    \PY@tok{#1}\expandafter\PY@toks\fi}
\def\PY@do#1{\PY@bc{\PY@tc{\PY@ul{%
    \PY@it{\PY@bf{\PY@ff{#1}}}}}}}
\def\PY#1#2{\PY@reset\PY@toks#1+\relax+\PY@do{#2}}

\@namedef{PY@tok@w}{\def\PY@tc##1{\textcolor[rgb]{0.73,0.73,0.73}{##1}}}
\@namedef{PY@tok@c}{\let\PY@it=\textit\def\PY@tc##1{\textcolor[rgb]{0.24,0.48,0.48}{##1}}}
\@namedef{PY@tok@cp}{\def\PY@tc##1{\textcolor[rgb]{0.61,0.40,0.00}{##1}}}
\@namedef{PY@tok@k}{\let\PY@bf=\textbf\def\PY@tc##1{\textcolor[rgb]{0.00,0.50,0.00}{##1}}}
\@namedef{PY@tok@kp}{\def\PY@tc##1{\textcolor[rgb]{0.00,0.50,0.00}{##1}}}
\@namedef{PY@tok@kt}{\def\PY@tc##1{\textcolor[rgb]{0.69,0.00,0.25}{##1}}}
\@namedef{PY@tok@o}{\def\PY@tc##1{\textcolor[rgb]{0.40,0.40,0.40}{##1}}}
\@namedef{PY@tok@ow}{\let\PY@bf=\textbf\def\PY@tc##1{\textcolor[rgb]{0.67,0.13,1.00}{##1}}}
\@namedef{PY@tok@nb}{\def\PY@tc##1{\textcolor[rgb]{0.00,0.50,0.00}{##1}}}
\@namedef{PY@tok@nf}{\def\PY@tc##1{\textcolor[rgb]{0.00,0.00,1.00}{##1}}}
\@namedef{PY@tok@nc}{\let\PY@bf=\textbf\def\PY@tc##1{\textcolor[rgb]{0.00,0.00,1.00}{##1}}}
\@namedef{PY@tok@nn}{\let\PY@bf=\textbf\def\PY@tc##1{\textcolor[rgb]{0.00,0.00,1.00}{##1}}}
\@namedef{PY@tok@ne}{\let\PY@bf=\textbf\def\PY@tc##1{\textcolor[rgb]{0.80,0.25,0.22}{##1}}}
\@namedef{PY@tok@nv}{\def\PY@tc##1{\textcolor[rgb]{0.10,0.09,0.49}{##1}}}
\@namedef{PY@tok@no}{\def\PY@tc##1{\textcolor[rgb]{0.53,0.00,0.00}{##1}}}
\@namedef{PY@tok@nl}{\def\PY@tc##1{\textcolor[rgb]{0.46,0.46,0.00}{##1}}}
\@namedef{PY@tok@ni}{\let\PY@bf=\textbf\def\PY@tc##1{\textcolor[rgb]{0.44,0.44,0.44}{##1}}}
\@namedef{PY@tok@na}{\def\PY@tc##1{\textcolor[rgb]{0.41,0.47,0.13}{##1}}}
\@namedef{PY@tok@nt}{\let\PY@bf=\textbf\def\PY@tc##1{\textcolor[rgb]{0.00,0.50,0.00}{##1}}}
\@namedef{PY@tok@nd}{\def\PY@tc##1{\textcolor[rgb]{0.67,0.13,1.00}{##1}}}
\@namedef{PY@tok@s}{\def\PY@tc##1{\textcolor[rgb]{0.73,0.13,0.13}{##1}}}
\@namedef{PY@tok@sd}{\let\PY@it=\textit\def\PY@tc##1{\textcolor[rgb]{0.73,0.13,0.13}{##1}}}
\@namedef{PY@tok@si}{\let\PY@bf=\textbf\def\PY@tc##1{\textcolor[rgb]{0.64,0.35,0.47}{##1}}}
\@namedef{PY@tok@se}{\let\PY@bf=\textbf\def\PY@tc##1{\textcolor[rgb]{0.67,0.36,0.12}{##1}}}
\@namedef{PY@tok@sr}{\def\PY@tc##1{\textcolor[rgb]{0.64,0.35,0.47}{##1}}}
\@namedef{PY@tok@ss}{\def\PY@tc##1{\textcolor[rgb]{0.10,0.09,0.49}{##1}}}
\@namedef{PY@tok@sx}{\def\PY@tc##1{\textcolor[rgb]{0.00,0.50,0.00}{##1}}}
\@namedef{PY@tok@m}{\def\PY@tc##1{\textcolor[rgb]{0.40,0.40,0.40}{##1}}}
\@namedef{PY@tok@gh}{\let\PY@bf=\textbf\def\PY@tc##1{\textcolor[rgb]{0.00,0.00,0.50}{##1}}}
\@namedef{PY@tok@gu}{\let\PY@bf=\textbf\def\PY@tc##1{\textcolor[rgb]{0.50,0.00,0.50}{##1}}}
\@namedef{PY@tok@gd}{\def\PY@tc##1{\textcolor[rgb]{0.63,0.00,0.00}{##1}}}
\@namedef{PY@tok@gi}{\def\PY@tc##1{\textcolor[rgb]{0.00,0.52,0.00}{##1}}}
\@namedef{PY@tok@gr}{\def\PY@tc##1{\textcolor[rgb]{0.89,0.00,0.00}{##1}}}
\@namedef{PY@tok@ge}{\let\PY@it=\textit}
\@namedef{PY@tok@gs}{\let\PY@bf=\textbf}
\@namedef{PY@tok@ges}{\let\PY@bf=\textbf\let\PY@it=\textit}
\@namedef{PY@tok@gp}{\let\PY@bf=\textbf\def\PY@tc##1{\textcolor[rgb]{0.00,0.00,0.50}{##1}}}
\@namedef{PY@tok@go}{\def\PY@tc##1{\textcolor[rgb]{0.44,0.44,0.44}{##1}}}
\@namedef{PY@tok@gt}{\def\PY@tc##1{\textcolor[rgb]{0.00,0.27,0.87}{##1}}}
\@namedef{PY@tok@err}{\def\PY@bc##1{{\setlength{\fboxsep}{\string -\fboxrule}\fcolorbox[rgb]{1.00,0.00,0.00}{1,1,1}{\strut ##1}}}}
\@namedef{PY@tok@kc}{\let\PY@bf=\textbf\def\PY@tc##1{\textcolor[rgb]{0.00,0.50,0.00}{##1}}}
\@namedef{PY@tok@kd}{\let\PY@bf=\textbf\def\PY@tc##1{\textcolor[rgb]{0.00,0.50,0.00}{##1}}}
\@namedef{PY@tok@kn}{\let\PY@bf=\textbf\def\PY@tc##1{\textcolor[rgb]{0.00,0.50,0.00}{##1}}}
\@namedef{PY@tok@kr}{\let\PY@bf=\textbf\def\PY@tc##1{\textcolor[rgb]{0.00,0.50,0.00}{##1}}}
\@namedef{PY@tok@bp}{\def\PY@tc##1{\textcolor[rgb]{0.00,0.50,0.00}{##1}}}
\@namedef{PY@tok@fm}{\def\PY@tc##1{\textcolor[rgb]{0.00,0.00,1.00}{##1}}}
\@namedef{PY@tok@vc}{\def\PY@tc##1{\textcolor[rgb]{0.10,0.09,0.49}{##1}}}
\@namedef{PY@tok@vg}{\def\PY@tc##1{\textcolor[rgb]{0.10,0.09,0.49}{##1}}}
\@namedef{PY@tok@vi}{\def\PY@tc##1{\textcolor[rgb]{0.10,0.09,0.49}{##1}}}
\@namedef{PY@tok@vm}{\def\PY@tc##1{\textcolor[rgb]{0.10,0.09,0.49}{##1}}}
\@namedef{PY@tok@sa}{\def\PY@tc##1{\textcolor[rgb]{0.73,0.13,0.13}{##1}}}
\@namedef{PY@tok@sb}{\def\PY@tc##1{\textcolor[rgb]{0.73,0.13,0.13}{##1}}}
\@namedef{PY@tok@sc}{\def\PY@tc##1{\textcolor[rgb]{0.73,0.13,0.13}{##1}}}
\@namedef{PY@tok@dl}{\def\PY@tc##1{\textcolor[rgb]{0.73,0.13,0.13}{##1}}}
\@namedef{PY@tok@s2}{\def\PY@tc##1{\textcolor[rgb]{0.73,0.13,0.13}{##1}}}
\@namedef{PY@tok@sh}{\def\PY@tc##1{\textcolor[rgb]{0.73,0.13,0.13}{##1}}}
\@namedef{PY@tok@s1}{\def\PY@tc##1{\textcolor[rgb]{0.73,0.13,0.13}{##1}}}
\@namedef{PY@tok@mb}{\def\PY@tc##1{\textcolor[rgb]{0.40,0.40,0.40}{##1}}}
\@namedef{PY@tok@mf}{\def\PY@tc##1{\textcolor[rgb]{0.40,0.40,0.40}{##1}}}
\@namedef{PY@tok@mh}{\def\PY@tc##1{\textcolor[rgb]{0.40,0.40,0.40}{##1}}}
\@namedef{PY@tok@mi}{\def\PY@tc##1{\textcolor[rgb]{0.40,0.40,0.40}{##1}}}
\@namedef{PY@tok@il}{\def\PY@tc##1{\textcolor[rgb]{0.40,0.40,0.40}{##1}}}
\@namedef{PY@tok@mo}{\def\PY@tc##1{\textcolor[rgb]{0.40,0.40,0.40}{##1}}}
\@namedef{PY@tok@ch}{\let\PY@it=\textit\def\PY@tc##1{\textcolor[rgb]{0.24,0.48,0.48}{##1}}}
\@namedef{PY@tok@cm}{\let\PY@it=\textit\def\PY@tc##1{\textcolor[rgb]{0.24,0.48,0.48}{##1}}}
\@namedef{PY@tok@cpf}{\let\PY@it=\textit\def\PY@tc##1{\textcolor[rgb]{0.24,0.48,0.48}{##1}}}
\@namedef{PY@tok@c1}{\let\PY@it=\textit\def\PY@tc##1{\textcolor[rgb]{0.24,0.48,0.48}{##1}}}
\@namedef{PY@tok@cs}{\let\PY@it=\textit\def\PY@tc##1{\textcolor[rgb]{0.24,0.48,0.48}{##1}}}

\def\PYZbs{\char`\\}
\def\PYZus{\char`\_}
\def\PYZob{\char`\{}
\def\PYZcb{\char`\}}
\def\PYZca{\char`\^}
\def\PYZam{\char`\&}
\def\PYZlt{\char`\<}
\def\PYZgt{\char`\>}
\def\PYZsh{\char`\#}
\def\PYZpc{\char`\%}
\def\PYZdl{\char`\$}
\def\PYZhy{\char`\-}
\def\PYZsq{\char`\'}
\def\PYZdq{\char`\"}
\def\PYZti{\char`\~}
% for compatibility with earlier versions
\def\PYZat{@}
\def\PYZlb{[}
\def\PYZrb{]}
\makeatother


    % For linebreaks inside Verbatim environment from package fancyvrb.
    \makeatletter
        \newbox\Wrappedcontinuationbox
        \newbox\Wrappedvisiblespacebox
        \newcommand*\Wrappedvisiblespace {\textcolor{red}{\textvisiblespace}}
        \newcommand*\Wrappedcontinuationsymbol {\textcolor{red}{\llap{\tiny$\m@th\hookrightarrow$}}}
        \newcommand*\Wrappedcontinuationindent {3ex }
        \newcommand*\Wrappedafterbreak {\kern\Wrappedcontinuationindent\copy\Wrappedcontinuationbox}
        % Take advantage of the already applied Pygments mark-up to insert
        % potential linebreaks for TeX processing.
        %        {, <, #, %, $, ' and ": go to next line.
        %        _, }, ^, &, >, - and ~: stay at end of broken line.
        % Use of \textquotesingle for straight quote.
        \newcommand*\Wrappedbreaksatspecials {%
            \def\PYGZus{\discretionary{\char`\_}{\Wrappedafterbreak}{\char`\_}}%
            \def\PYGZob{\discretionary{}{\Wrappedafterbreak\char`\{}{\char`\{}}%
            \def\PYGZcb{\discretionary{\char`\}}{\Wrappedafterbreak}{\char`\}}}%
            \def\PYGZca{\discretionary{\char`\^}{\Wrappedafterbreak}{\char`\^}}%
            \def\PYGZam{\discretionary{\char`\&}{\Wrappedafterbreak}{\char`\&}}%
            \def\PYGZlt{\discretionary{}{\Wrappedafterbreak\char`\<}{\char`\<}}%
            \def\PYGZgt{\discretionary{\char`\>}{\Wrappedafterbreak}{\char`\>}}%
            \def\PYGZsh{\discretionary{}{\Wrappedafterbreak\char`\#}{\char`\#}}%
            \def\PYGZpc{\discretionary{}{\Wrappedafterbreak\char`\%}{\char`\%}}%
            \def\PYGZdl{\discretionary{}{\Wrappedafterbreak\char`\$}{\char`\$}}%
            \def\PYGZhy{\discretionary{\char`\-}{\Wrappedafterbreak}{\char`\-}}%
            \def\PYGZsq{\discretionary{}{\Wrappedafterbreak\textquotesingle}{\textquotesingle}}%
            \def\PYGZdq{\discretionary{}{\Wrappedafterbreak\char`\"}{\char`\"}}%
            \def\PYGZti{\discretionary{\char`\~}{\Wrappedafterbreak}{\char`\~}}%
        }
        % Some characters . , ; ? ! / are not pygmentized.
        % This macro makes them "active" and they will insert potential linebreaks
        \newcommand*\Wrappedbreaksatpunct {%
            \lccode`\~`\.\lowercase{\def~}{\discretionary{\hbox{\char`\.}}{\Wrappedafterbreak}{\hbox{\char`\.}}}%
            \lccode`\~`\,\lowercase{\def~}{\discretionary{\hbox{\char`\,}}{\Wrappedafterbreak}{\hbox{\char`\,}}}%
            \lccode`\~`\;\lowercase{\def~}{\discretionary{\hbox{\char`\;}}{\Wrappedafterbreak}{\hbox{\char`\;}}}%
            \lccode`\~`\:\lowercase{\def~}{\discretionary{\hbox{\char`\:}}{\Wrappedafterbreak}{\hbox{\char`\:}}}%
            \lccode`\~`\?\lowercase{\def~}{\discretionary{\hbox{\char`\?}}{\Wrappedafterbreak}{\hbox{\char`\?}}}%
            \lccode`\~`\!\lowercase{\def~}{\discretionary{\hbox{\char`\!}}{\Wrappedafterbreak}{\hbox{\char`\!}}}%
            \lccode`\~`\/\lowercase{\def~}{\discretionary{\hbox{\char`\/}}{\Wrappedafterbreak}{\hbox{\char`\/}}}%
            \catcode`\.\active
            \catcode`\,\active
            \catcode`\;\active
            \catcode`\:\active
            \catcode`\?\active
            \catcode`\!\active
            \catcode`\/\active
            \lccode`\~`\~
        }
    \makeatother

    \let\OriginalVerbatim=\Verbatim
    \makeatletter
    \renewcommand{\Verbatim}[1][1]{%
        %\parskip\z@skip
        \sbox\Wrappedcontinuationbox {\Wrappedcontinuationsymbol}%
        \sbox\Wrappedvisiblespacebox {\FV@SetupFont\Wrappedvisiblespace}%
        \def\FancyVerbFormatLine ##1{\hsize\linewidth
            \vtop{\raggedright\hyphenpenalty\z@\exhyphenpenalty\z@
                \doublehyphendemerits\z@\finalhyphendemerits\z@
                \strut ##1\strut}%
        }%
        % If the linebreak is at a space, the latter will be displayed as visible
        % space at end of first line, and a continuation symbol starts next line.
        % Stretch/shrink are however usually zero for typewriter font.
        \def\FV@Space {%
            \nobreak\hskip\z@ plus\fontdimen3\font minus\fontdimen4\font
            \discretionary{\copy\Wrappedvisiblespacebox}{\Wrappedafterbreak}
            {\kern\fontdimen2\font}%
        }%

        % Allow breaks at special characters using \PYG... macros.
        \Wrappedbreaksatspecials
        % Breaks at punctuation characters . , ; ? ! and / need catcode=\active
        \OriginalVerbatim[#1,codes*=\Wrappedbreaksatpunct]%
    }
    \makeatother

    % Exact colors from NB
    \definecolor{incolor}{HTML}{303F9F}
    \definecolor{outcolor}{HTML}{D84315}
    \definecolor{cellborder}{HTML}{CFCFCF}
    \definecolor{cellbackground}{HTML}{F7F7F7}

    % prompt
    \makeatletter
    \newcommand{\boxspacing}{\kern\kvtcb@left@rule\kern\kvtcb@boxsep}
    \makeatother
    \newcommand{\prompt}[4]{
        {\ttfamily\llap{{\color{#2}[#3]:\hspace{3pt}#4}}\vspace{-\baselineskip}}
    }
    

    
    % Prevent overflowing lines due to hard-to-break entities
    \sloppy
    % Setup hyperref package
    \hypersetup{
      breaklinks=true,  % so long urls are correctly broken across lines
      colorlinks=true,
      urlcolor=urlcolor,
      linkcolor=linkcolor,
      citecolor=citecolor,
      }
    % Slightly bigger margins than the latex defaults
    
    \geometry{verbose,tmargin=1in,bmargin=1in,lmargin=1in,rmargin=1in}
    
    

\begin{document}
    
    \maketitle
    
    

    
    \section{Python
数据数学分析核心课件}\label{python-ux6570ux636eux6570ux5b66ux5206ux6790ux6838ux5fc3ux8bfeux4ef6}

\subsection{目录}\label{ux76eeux5f55}

\subsection{1.
引言:从工程代码回归数据本质}\label{ux5f15ux8a00ux4eceux5de5ux7a0bux4ee3ux7801ux56deux5f52ux6570ux636eux672cux8d28}

\subsubsection{1.1 为什么是
Python？}\label{ux4e3aux4ec0ux4e48ux662f-python}

在传统的软件后端开发中, 关注点通常是 \textbf{状态管理、网络
I/O、并发安全与面向对象的解耦抽象}. 然而,
当进入数据分析与科学计算领域时, 思维模型需要发生根本性转变:

\begin{verbatim}
[传统后端工程思维]                     [现代数据分析/矩阵思维]
对象封装(OOP)                   ->   数据与操作分离(Data-Oriented Design)
逐个遍历(for item in list)       ->   批量向量化计算(SIMD / Tensor Processing)
复杂控制流(if-else / 状态机)      ->   布尔掩码与索引选择(Mask & Fancy Indexing)
显式多线程/锁竞争                 ->   底层 C/Fortran/BLAS 内存连续并行计算
\end{verbatim}

Python 之所以能够统治现代数据科学与 AI 领域,
并非因为其解释器运行效率高(相反, CPython 的解释器开销与 GIL
限制众所周知), 而是因为它的 \textbf{``胶水特性'' 与 C-API 扩展能力}:

\begin{enumerate}
\def\labelenumi{\arabic{enumi}.}
\tightlist
\item
  上层提供极具表现力的动态语法和交互式 REPL 环境;
\item
  下层无缝对接 OpenBLAS、MKL、CUDA 等底层硬件加速库;
\item
  形成了从数据摄取(Pandas/Arrow)、处理(NumPy/SciPy)、可视化(Matplotlib/Seaborn)到模型训练(PyTorch/Scikit-learn)的完整闭环.
\end{enumerate}

\begin{quote}
为什么 Python 明明什么都能算, 我们还需要 NumPy 和 Pandas？
\end{quote}

\begin{quote}
NumPy 和 Pandas 真正带来的价值, 是``多了一堆 API'',
还是``改变了数据和计算的抽象方式''？
\end{quote}

\subsubsection{我们真正要学的不是
API}\label{ux6211ux4eecux771fux6b63ux8981ux5b66ux7684ux4e0dux662f-api}

假设现在拿到 200 万条出租车订单.

我们需要分析:

\begin{Shaded}
\begin{Highlighting}[]
\NormalTok{哪些订单不可信？}

\NormalTok{什么时候最忙？}

\NormalTok{什么时候流水最高？}

\NormalTok{哪些区域最热门？}

\NormalTok{哪些订单看起来异常？}

\NormalTok{为什么同一个计算, 有人的代码跑几十秒, 有人的代码只跑几秒？}

\NormalTok{怎么把这些分析变成公司里以后还能复用的工具？}
\end{Highlighting}
\end{Shaded}

当然, 可以全部使用 Python:

\begin{Shaded}
\begin{Highlighting}[]
\ControlFlowTok{for}\NormalTok{ row }\KeywordTok{in}\NormalTok{ rows:}
\NormalTok{    ...}
\end{Highlighting}
\end{Shaded}

但问题并不在于 Python \textbf{能不能做}.

真正的问题是:

\begin{quote}
\textbf{Python 原生的数据结构和逐对象执行模型,
并不是专门针对百万级同构数值计算和二维表格分析设计的. }
\end{quote}

但是 Numpy 和 Pandas 是 Python 中专门为数据分析设计的

\begin{enumerate}
\def\labelenumi{\arabic{enumi}.}
\tightlist
\item
  NumPy 的核心是同构 \(N\) 维数组 \texttt{ndarray}；官方文档将其定义为
  NumPy 的核心多维数组结构,
  并围绕数组提供高效的数学、逻辑、选择、排序、线性代数等运算.
\item
  Pandas 则在数组之上增加了
  \textbf{标签、索引、缺失值、异构列、自动对齐、分组、连接和时间序列}
  等更贴近表格数据分析的抽象；\texttt{Index} 本身就是 Pandas
  用于索引和对齐的轴标签对象.
\end{enumerate}

\begin{Shaded}
\begin{Highlighting}[]
\NormalTok{flowchart TD}
\NormalTok{    A["Python 原生对象\textless{}br/\textgreater{}list / dict / Python object"] {-}{-}\textgreater{} B["NumPy ndarray"]}
\NormalTok{    B {-}{-}\textgreater{} C["Pandas Series / DataFrame"]}
\NormalTok{    C {-}{-}\textgreater{} D["业务分析"]}
\NormalTok{    D {-}{-}\textgreater{} E["性能工程"]}
\NormalTok{    E {-}{-}\textgreater{} F["领域工具 / Pipeline / Accessor"]}

\NormalTok{    B {-}{-}{-} B1["dtype / shape / strides"]}
\NormalTok{    B {-}{-}{-} B2["ufunc / vectorization"]}
\NormalTok{    B {-}{-}{-} B3["broadcasting"]}

\NormalTok{    C {-}{-}{-} C1["Index / alignment"]}
\NormalTok{    C {-}{-}{-} C2["groupby / merge"]}
\NormalTok{    C {-}{-}{-} C3["resample / rolling"]}
\end{Highlighting}
\end{Shaded}

一句话概括: \textgreater{} \textbf{NumPy
解决``如何高效计算数组''；Pandas
解决``如何带着业务语义高效组织、查询、组合和分析表格数据''. }

\subsubsection{出租车数据集简介}\label{ux51faux79dfux8f66ux6570ux636eux96c6ux7b80ux4ecb}

NYC TLC Yellow Taxi Trip Records 是纽约市出租车与豪华轿车委员会(NYC Taxi
\& Limousine Commission, TLC)公开发布的纽约市黄色出租车行程记录数据.

\subsubsection{贯穿本节的七个问题}\label{ux8d2fux7a7fux672cux8282ux7684ux4e03ux4e2aux95eeux9898}

{\def\LTcaptype{none} % do not increment counter
\begin{longtable}[]{@{}
  >{\raggedright\arraybackslash}p{(\linewidth - 4\tabcolsep) * \real{0.3333}}
  >{\raggedright\arraybackslash}p{(\linewidth - 4\tabcolsep) * \real{0.3333}}
  >{\raggedright\arraybackslash}p{(\linewidth - 4\tabcolsep) * \real{0.3333}}@{}}
\toprule\noalign{}
\begin{minipage}[b]{\linewidth}\raggedright
问题
\end{minipage} & \begin{minipage}[b]{\linewidth}\raggedright
业务问题
\end{minipage} & \begin{minipage}[b]{\linewidth}\raggedright
逐步引出的技术
\end{minipage} \\
\midrule\noalign{}
\endhead
\bottomrule\noalign{}
\endlastfoot
\textbf{Q1} & 这 200 万条数据可靠吗？ &
\texttt{dtype}、缺失值、mask、\texttt{loc}、\texttt{query} \\
\textbf{Q2} & 一天中什么时候最忙？ &
\texttt{datetime}、\texttt{groupby}、\texttt{agg} \\
\textbf{Q3} & 什么时间段的流水和单位运营时间收入最高？ &
向量化、\texttt{assign}、\texttt{groupby}、派生指标 \\
\textbf{Q4} & 哪些上下车区域最热门？ &
\texttt{value\_counts}、\texttt{merge}、连接校验 \\
\textbf{Q5} & 哪些订单和时间窗口可能异常？ &
\texttt{where}、\texttt{select}、\texttt{quantile}、\texttt{transform}、\texttt{resample}、\texttt{rolling} \\
\textbf{Q6} & 为什么同一个分析可以相差一个数量级甚至更多？ &
\texttt{Python\ loop}、\texttt{apply}、\texttt{NumPy}、\texttt{Numba}、\texttt{eval/query}、内存优化 \\
\textbf{Q7} & 如何把分析脚本变成公司可以重复使用的工具？ &
\texttt{pipe}、纯函数、\texttt{Accessor}、领域 API \\
\end{longtable}
}

\subsubsection{NumPy / Pandas
为什么存在}\label{numpy-pandas-ux4e3aux4ec0ux4e48ux5b58ux5728}

\paragraph{问题}\label{ux95eeux9898}

假设我们有 200 万个金额:

\begin{Shaded}
\begin{Highlighting}[]
\NormalTok{prices }\OperatorTok{=}\NormalTok{ [...]}
\end{Highlighting}
\end{Shaded}

需要:

\begin{Shaded}
\begin{Highlighting}[]
\NormalTok{全部上涨 5\%}
\NormalTok{再增加 3 元服务费}
\NormalTok{过滤异常价格}
\NormalTok{求平均值}
\end{Highlighting}
\end{Shaded}

\paragraph{直接使用 Python List
实现}\label{ux76f4ux63a5ux4f7fux7528-python-list-ux5b9eux73b0}

\begin{Shaded}
\begin{Highlighting}[]
\CommentTok{\# 直接使用 List 存储与计算}
\NormalTok{result }\OperatorTok{=}\NormalTok{ []}

\ControlFlowTok{for}\NormalTok{ price }\KeywordTok{in}\NormalTok{ prices:}
\NormalTok{    result.append(price }\OperatorTok{*} \FloatTok{1.05} \OperatorTok{+} \DecValTok{3}\NormalTok{)}
\end{Highlighting}
\end{Shaded}

问题是我们把 \textbf{200 万次循环控制} 交给了 Python 层. 具体来说,
Python 原生 \texttt{List} 在这种大批量数值计算场景下存在两大明显短板:

\textbf{1. 存储方面}

Python \texttt{List} 并不是一块连续的数值缓冲区, 而是一个
\textbf{对象引用数组}. 它存储的每一个元素都是一个独立的 Python
对象(这里是 \texttt{float}), 每个对象都带有引用计数、类型指针等元信息.
以 200 万个金额为例:

\begin{Shaded}
\begin{Highlighting}[]
\NormalTok{prices }\OperatorTok{=}\NormalTok{ [}\FloatTok{42.5}\NormalTok{, }\FloatTok{17.3}\NormalTok{, }\FloatTok{58.9}\NormalTok{, ...]   }\CommentTok{\# 200 万个 float 对象}
\end{Highlighting}
\end{Shaded}

\begin{itemize}
\tightlist
\item
  \texttt{List} 本身只存 200 万个\textbf{指针}, 每个指针 8 字节(64
  位系统)；
\item
  每一个 \texttt{float} 对象本身还要额外占用约 24 字节(CPython 对象头)；
\item
  于是 \texttt{List} 存储 200 万个浮点数的实际内存, 大约是相同规模
  \texttt{ndarray} 的 \textbf{3\textasciitilde5 倍}.
\end{itemize}

\begin{Shaded}
\begin{Highlighting}[]
\NormalTok{List:}
\NormalTok{  指针数组(8B × 200万) + 每个 float 对象(24B × 200万)}
\NormalTok{  ≈ 16 MB + 48 MB = 64 MB 左右}

\NormalTok{ndarray(float64):}
\NormalTok{  连续缓冲区(8B × 200万)}
\NormalTok{  ≈ 16 MB}
\end{Highlighting}
\end{Shaded}

内存翻倍, 还意味着缓存命中率更低、CPU 需要跨越更多随机内存地址.

\textbf{2. 计算方面}

\begin{Shaded}
\begin{Highlighting}[]
\NormalTok{result }\OperatorTok{=}\NormalTok{ []}
\ControlFlowTok{for}\NormalTok{ price }\KeywordTok{in}\NormalTok{ prices:}
\NormalTok{    result.append(price }\OperatorTok{*} \FloatTok{1.05} \OperatorTok{+} \DecValTok{3}\NormalTok{)}
\end{Highlighting}
\end{Shaded}

这段循环的执行, 每一步都发生在 \textbf{Python 解释器层},
而不是底层机器码层:

\begin{itemize}
\tightlist
\item
  每次 \texttt{for} 迭代都要经历字节码分发(bytecode dispatch)；
\item
  每次读取 \texttt{price} 都要做一次 \textbf{拆箱}(unboxing, 从 Python
  对象还原成 C double)；
\item
  每次 \texttt{price\ *\ 1.05\ +\ 3} 结果要重新
  \textbf{装箱}(boxing)成新的 \texttt{float} 对象；
\item
  每次 \texttt{append}
  还要调用方法、检查容量、可能触发列表扩容与内存重分配.
\end{itemize}

当循环次数上升到 200 万, 这些\textbf{解释器级开销被放大 200 万倍},
成为无可忽视的性能瓶颈. 这也是为什么同样的计算, 纯 Python 循环往往比
NumPy 向量化写法慢一到两个数量级------\textbf{真正慢的不是''算'',
而是''每次都要让解释器来驱动''}.

\paragraph{使用 Numpy 和 Pandas
实现}\label{ux4f7fux7528-numpy-ux548c-pandas-ux5b9eux73b0}

NumPy:

\begin{Shaded}
\begin{Highlighting}[]
\NormalTok{result }\OperatorTok{=}\NormalTok{ prices }\OperatorTok{*} \FloatTok{1.05} \OperatorTok{+} \DecValTok{3}
\end{Highlighting}
\end{Shaded}

Pandas:

\begin{Shaded}
\begin{Highlighting}[]
\NormalTok{result }\OperatorTok{=}\NormalTok{ (}
\NormalTok{    df}
\NormalTok{    .query(}\StringTok{"fare\_amount \textgreater{} 0"}\NormalTok{)}
\NormalTok{    .groupby(}\StringTok{"pickup\_hour"}\NormalTok{)}
\NormalTok{    .agg(avg\_fare}\OperatorTok{=}\NormalTok{(}\StringTok{"fare\_amount"}\NormalTok{, }\StringTok{"mean"}\NormalTok{))}
\NormalTok{)}
\end{Highlighting}
\end{Shaded}

\subsubsection{Numpy 和 Pandas
的特点}\label{numpy-ux548c-pandas-ux7684ux7279ux70b9}

\textbf{NumPy 解决的痛点(纯数值与张量计算):} 在科学计算场景下, 原生
Python 的 \texttt{list}
存储机制会导致严重的内存碎片化(频繁指针解引用)和缓存未命中(Cache Miss),
同时动态类型解释器导致循环极慢.NumPy
的诞生终结了这五大痛点:解决内存空间碎片化、消除解释器动态类型开销、填补多维张量代数的空白、终结切片的高昂拷贝代价,
并打通底层 C/BLAS 和 SIMD 硬件加速生态.

\textbf{Pandas 解决的痛点(异构业务数据与关系代数):}
真实世界的业务数据不是纯数学矩阵,
而是包含字符串、时间戳、布尔值的异构表格, 常存在缺失值和错位情况.Pandas
解决了: 1. 如何容纳异构数据类型(字符串、数值混合). 2.
在连接表或对齐时如何依据业务``键(Label)''而非纯``位置(Index)''进行运算.
3. 如何像 SQL 一样使用 GroupBy、Join 等关系代数算子. 4.
如何优雅且不中断程序地处理 NaN 缺失数据.

\paragraph{核心设计哲学}\label{ux6838ux5fc3ux8bbeux8ba1ux54f2ux5b66}

\subparagraph{\texorpdfstring{NumPy
的灵魂核心:\texttt{ndarray}(N-dimensional
array)}{NumPy 的灵魂核心:ndarray(N-dimensional array)}}\label{numpy-ux7684ux7075ux9b42ux6838ux5fc3ndarrayn-dimensional-array}

\texttt{ndarray}(N-Dimensional Array)
的本质是``同构多维连续内存缓冲区的视图抽象(N-dimensional Array /
Tensor)'', 矩阵(Matrix)是它在 2
维空间下的一个特例.它的设计哲学是:\textbf{在动态语言中实现静态语言级别的内存布局.}

\begin{itemize}
\tightlist
\item
  \textbf{元数据解耦与视图(View):} \texttt{ndarray} 将数据的
  \texttt{Header}(包含 shape、dtype、strides)与底层的
  \texttt{Data\ Buffer}(连续物理内存)彻底分离.这使得切片、转置、Reshape
  都只需要极小开销修改 Header 且不需要复制物理内存.
\item
  \textbf{齐次强类型与 SIMD:} 数据块内类型必须一致, 使得 CPU
  能够跨步寻址(索引 × 大小), 并直接投喂给硬件向量化指令.
\item
  \textbf{面向数组编程:} 把底层长循环隐式推给 C 语言层面,
  上层仅做``张量形态''的声明式表达.
\end{itemize}

\subparagraph{\texorpdfstring{Pandas 的根本基石:\texttt{Series} 和
\texttt{DataFrame}}{Pandas 的根本基石:Series 和 DataFrame}}\label{pandas-ux7684ux6839ux672cux57faux77f3series-ux548c-dataframe}

如果 NumPy 是``同质数学张量'', Pandas 则是``带标签的异构列式数据框''. *
\textbf{列式存储架构:} DataFrame 本质是多个一维 \texttt{Series}
的字典集合.在底层,
同类型的列会被组织为内存连续的块.列式存储使得``求整列均值''等分析操作在物理上读取连续内存,
极其高效. * \textbf{标签驱动(Label-Based Alignment):} 与 NumPy
严格的矩阵位置形状对齐不同, Series 和 DataFrame
的核心哲学在于``业务主键(Index)绑定''.两张表相加时, Pandas
会自动依据相同的 Index 标签进行匹配(如果没有则视为 NaN),
哪怕它们数据行数或顺序完全不同.这规避了业务分析中最致命的错位问题.

    \subsection{4.
课程环境与初始化}\label{ux8bfeux7a0bux73afux5883ux4e0eux521dux59cbux5316}

\subsubsection{Numpy 与 Pandas
的版本}\label{numpy-ux4e0e-pandas-ux7684ux7248ux672c}

本课程 Numpy 与 Pandas 需要安装与配套的版本如下:

Package Version, matplotlib 3.9.4, numba 0.60.0, numpy 2.0.2, pandas
2.3.3, pyarrow 21.0.0, jupyterlab 4.5.10

\subsubsection{Jupyter Notebook
初始化}\label{jupyter-notebook-ux521dux59cbux5316}

    \begin{tcolorbox}[breakable, size=fbox, boxrule=1pt, pad at break*=1mm,colback=cellbackground, colframe=cellborder]
\prompt{In}{incolor}{1}{\boxspacing}
\begin{Verbatim}[commandchars=\\\{\}]
\PY{k+kn}{from}\PY{+w}{ }\PY{n+nn}{\PYZus{}\PYZus{}future\PYZus{}\PYZus{}}\PY{+w}{ }\PY{k+kn}{import} \PY{n}{annotations}

\PY{k+kn}{import}\PY{+w}{ }\PY{n+nn}{gc}
\PY{k+kn}{import}\PY{+w}{ }\PY{n+nn}{math}
\PY{k+kn}{import}\PY{+w}{ }\PY{n+nn}{os}
\PY{k+kn}{import}\PY{+w}{ }\PY{n+nn}{platform}
\PY{k+kn}{import}\PY{+w}{ }\PY{n+nn}{sys}
\PY{k+kn}{import}\PY{+w}{ }\PY{n+nn}{time}
\PY{k+kn}{import}\PY{+w}{ }\PY{n+nn}{timeit}
\PY{k+kn}{from}\PY{+w}{ }\PY{n+nn}{pathlib}\PY{+w}{ }\PY{k+kn}{import} \PY{n}{Path}    
\PY{k+kn}{from}\PY{+w}{ }\PY{n+nn}{urllib}\PY{n+nn}{.}\PY{n+nn}{request}\PY{+w}{ }\PY{k+kn}{import} \PY{n}{urlretrieve}

\PY{k+kn}{import}\PY{+w}{ }\PY{n+nn}{matplotlib}\PY{n+nn}{.}\PY{n+nn}{pyplot}\PY{+w}{ }\PY{k}{as}\PY{+w}{ }\PY{n+nn}{plt}
\PY{k+kn}{import}\PY{+w}{ }\PY{n+nn}{numba}
\PY{k+kn}{import}\PY{+w}{ }\PY{n+nn}{numpy}\PY{+w}{ }\PY{k}{as}\PY{+w}{ }\PY{n+nn}{np}
\PY{k+kn}{import}\PY{+w}{ }\PY{n+nn}{pandas}\PY{+w}{ }\PY{k}{as}\PY{+w}{ }\PY{n+nn}{pd}
\PY{k+kn}{import}\PY{+w}{ }\PY{n+nn}{pyarrow}\PY{+w}{ }\PY{k}{as}\PY{+w}{ }\PY{n+nn}{pa}
\PY{k+kn}{import}\PY{+w}{ }\PY{n+nn}{pyarrow}\PY{n+nn}{.}\PY{n+nn}{parquet}\PY{+w}{ }\PY{k}{as}\PY{+w}{ }\PY{n+nn}{pq}
\PY{k+kn}{from}\PY{+w}{ }\PY{n+nn}{numba}\PY{+w}{ }\PY{k+kn}{import} \PY{n}{njit}

\PY{n+nb}{print}\PY{p}{(}\PY{l+s+s2}{\PYZdq{}}\PY{l+s+s2}{Python :}\PY{l+s+s2}{\PYZdq{}}\PY{p}{,} \PY{n}{sys}\PY{o}{.}\PY{n}{version}\PY{o}{.}\PY{n}{split}\PY{p}{(}\PY{p}{)}\PY{p}{[}\PY{l+m+mi}{0}\PY{p}{]}\PY{p}{)}
\PY{n+nb}{print}\PY{p}{(}\PY{l+s+s2}{\PYZdq{}}\PY{l+s+s2}{NumPy  :}\PY{l+s+s2}{\PYZdq{}}\PY{p}{,} \PY{n}{np}\PY{o}{.}\PY{n}{\PYZus{}\PYZus{}version\PYZus{}\PYZus{}}\PY{p}{)}
\PY{n+nb}{print}\PY{p}{(}\PY{l+s+s2}{\PYZdq{}}\PY{l+s+s2}{Pandas :}\PY{l+s+s2}{\PYZdq{}}\PY{p}{,} \PY{n}{pd}\PY{o}{.}\PY{n}{\PYZus{}\PYZus{}version\PYZus{}\PYZus{}}\PY{p}{)}
\PY{n+nb}{print}\PY{p}{(}\PY{l+s+s2}{\PYZdq{}}\PY{l+s+s2}{PyArrow:}\PY{l+s+s2}{\PYZdq{}}\PY{p}{,} \PY{n}{pa}\PY{o}{.}\PY{n}{\PYZus{}\PYZus{}version\PYZus{}\PYZus{}}\PY{p}{)}
\PY{n+nb}{print}\PY{p}{(}\PY{l+s+s2}{\PYZdq{}}\PY{l+s+s2}{Numba  :}\PY{l+s+s2}{\PYZdq{}}\PY{p}{,} \PY{n}{numba}\PY{o}{.}\PY{n}{\PYZus{}\PYZus{}version\PYZus{}\PYZus{}}\PY{p}{)}
\PY{n+nb}{print}\PY{p}{(}\PY{l+s+s2}{\PYZdq{}}\PY{l+s+s2}{OS     :}\PY{l+s+s2}{\PYZdq{}}\PY{p}{,} \PY{n}{platform}\PY{o}{.}\PY{n}{platform}\PY{p}{(}\PY{p}{)}\PY{p}{)}

\PY{n}{pd}\PY{o}{.}\PY{n}{set\PYZus{}option}\PY{p}{(}\PY{l+s+s2}{\PYZdq{}}\PY{l+s+s2}{display.max\PYZus{}columns}\PY{l+s+s2}{\PYZdq{}}\PY{p}{,} \PY{l+m+mi}{100}\PY{p}{)}
\PY{n}{pd}\PY{o}{.}\PY{n}{set\PYZus{}option}\PY{p}{(}\PY{l+s+s2}{\PYZdq{}}\PY{l+s+s2}{display.float\PYZus{}format}\PY{l+s+s2}{\PYZdq{}}\PY{p}{,} \PY{k}{lambda} \PY{n}{x}\PY{p}{:} \PY{l+s+sa}{f}\PY{l+s+s2}{\PYZdq{}}\PY{l+s+si}{\PYZob{}}\PY{n}{x}\PY{l+s+si}{:}\PY{l+s+s2}{,.3f}\PY{l+s+si}{\PYZcb{}}\PY{l+s+s2}{\PYZdq{}}\PY{p}{)}

\PY{n}{DATA\PYZus{}DIR} \PY{o}{=} \PY{n}{Path}\PY{p}{(}\PY{l+s+s2}{\PYZdq{}}\PY{l+s+s2}{data}\PY{l+s+s2}{\PYZdq{}}\PY{p}{)}
\PY{n}{ARTIFACT\PYZus{}DIR} \PY{o}{=} \PY{n}{Path}\PY{p}{(}\PY{l+s+s2}{\PYZdq{}}\PY{l+s+s2}{artifacts}\PY{l+s+s2}{\PYZdq{}}\PY{p}{)}

\PY{n}{DATA\PYZus{}DIR}\PY{o}{.}\PY{n}{mkdir}\PY{p}{(}\PY{n}{exist\PYZus{}ok}\PY{o}{=}\PY{k+kc}{True}\PY{p}{)}
\PY{n}{ARTIFACT\PYZus{}DIR}\PY{o}{.}\PY{n}{mkdir}\PY{p}{(}\PY{n}{exist\PYZus{}ok}\PY{o}{=}\PY{k+kc}{True}\PY{p}{)}

\PY{n}{TAXI\PYZus{}PARQUET\PYZus{}PATH} \PY{o}{=} \PY{n}{DATA\PYZus{}DIR} \PY{o}{/} \PY{l+s+s2}{\PYZdq{}}\PY{l+s+s2}{yellow\PYZus{}tripdata\PYZus{}2025\PYZhy{}12.parquet}\PY{l+s+s2}{\PYZdq{}}

\PY{n}{TARGET\PYZus{}ROWS} \PY{o}{=} \PY{n+nb}{int}\PY{p}{(}\PY{n}{os}\PY{o}{.}\PY{n}{getenv}\PY{p}{(}\PY{l+s+s2}{\PYZdq{}}\PY{l+s+s2}{TAXI\PYZus{}ROWS}\PY{l+s+s2}{\PYZdq{}}\PY{p}{,} \PY{l+s+s2}{\PYZdq{}}\PY{l+s+s2}{2000000}\PY{l+s+s2}{\PYZdq{}}\PY{p}{)}\PY{p}{)}
\PY{n}{BENCH\PYZus{}ROWS} \PY{o}{=} \PY{n+nb}{int}\PY{p}{(}\PY{n}{os}\PY{o}{.}\PY{n}{getenv}\PY{p}{(}\PY{l+s+s2}{\PYZdq{}}\PY{l+s+s2}{TAXI\PYZus{}BENCH\PYZus{}ROWS}\PY{l+s+s2}{\PYZdq{}}\PY{p}{,} \PY{l+s+s2}{\PYZdq{}}\PY{l+s+s2}{100000}\PY{l+s+s2}{\PYZdq{}}\PY{p}{)}\PY{p}{)}
\PY{n}{IO\PYZus{}ROWS} \PY{o}{=} \PY{n+nb}{int}\PY{p}{(}\PY{n}{os}\PY{o}{.}\PY{n}{getenv}\PY{p}{(}\PY{l+s+s2}{\PYZdq{}}\PY{l+s+s2}{TAXI\PYZus{}IO\PYZus{}ROWS}\PY{l+s+s2}{\PYZdq{}}\PY{p}{,} \PY{n+nb}{str}\PY{p}{(}\PY{n}{TARGET\PYZus{}ROWS}\PY{p}{)}\PY{p}{)}\PY{p}{)}

\PY{n}{RANDOM\PYZus{}SEED} \PY{o}{=} \PY{l+m+mi}{42}

\PY{n+nb}{print}\PY{p}{(}\PY{l+s+sa}{f}\PY{l+s+s2}{\PYZdq{}}\PY{l+s+si}{\PYZob{}}\PY{n}{TARGET\PYZus{}ROWS}\PY{l+s+si}{=:}\PY{l+s+s2}{,}\PY{l+s+si}{\PYZcb{}}\PY{l+s+s2}{\PYZdq{}}\PY{p}{)}
\PY{n+nb}{print}\PY{p}{(}\PY{l+s+sa}{f}\PY{l+s+s2}{\PYZdq{}}\PY{l+s+si}{\PYZob{}}\PY{n}{BENCH\PYZus{}ROWS}\PY{l+s+si}{=:}\PY{l+s+s2}{,}\PY{l+s+si}{\PYZcb{}}\PY{l+s+s2}{\PYZdq{}}\PY{p}{)}
\PY{n+nb}{print}\PY{p}{(}\PY{l+s+sa}{f}\PY{l+s+s2}{\PYZdq{}}\PY{l+s+si}{\PYZob{}}\PY{n}{IO\PYZus{}ROWS}\PY{l+s+si}{=:}\PY{l+s+s2}{,}\PY{l+s+si}{\PYZcb{}}\PY{l+s+s2}{\PYZdq{}}\PY{p}{)}
\end{Verbatim}
\end{tcolorbox}

    \begin{Verbatim}[commandchars=\\\{\}]
Python : 3.9.25
NumPy  : 2.0.2
Pandas : 2.3.3
PyArrow: 21.0.0
Numba  : 0.60.0
OS     : macOS-15.7.7-arm64-arm-64bit
TARGET\_ROWS=2,000,000
BENCH\_ROWS=100,000
IO\_ROWS=2,000,000
    \end{Verbatim}

    \subsubsection{本课程统一 Benchmark
工具}\label{ux672cux8bfeux7a0bux7edfux4e00-benchmark-ux5de5ux5177}

    \begin{tcolorbox}[breakable, size=fbox, boxrule=1pt, pad at break*=1mm,colback=cellbackground, colframe=cellborder]
\prompt{In}{incolor}{2}{\boxspacing}
\begin{Verbatim}[commandchars=\\\{\}]
\PY{k}{def}\PY{+w}{ }\PY{n+nf}{benchmark}\PY{p}{(}
    \PY{n}{cases}\PY{p}{:} \PY{n+nb}{dict}\PY{p}{[}\PY{n+nb}{str}\PY{p}{,} \PY{n+nb}{callable}\PY{p}{]}\PY{p}{,}
    \PY{o}{*}\PY{p}{,}
    \PY{n}{repeat}\PY{p}{:} \PY{n+nb}{int} \PY{o}{=} \PY{l+m+mi}{5}\PY{p}{,}
    \PY{n}{number}\PY{p}{:} \PY{n+nb}{int} \PY{o}{=} \PY{l+m+mi}{1}\PY{p}{,}
    \PY{n}{warmup}\PY{p}{:} \PY{n+nb}{int} \PY{o}{=} \PY{l+m+mi}{1}\PY{p}{,}
\PY{p}{)} \PY{o}{\PYZhy{}}\PY{o}{\PYZgt{}} \PY{n}{pd}\PY{o}{.}\PY{n}{DataFrame}\PY{p}{:}
\PY{+w}{    }\PY{l+s+sd}{\PYZdq{}\PYZdq{}\PYZdq{}}
\PY{l+s+sd}{    对多个无参 callable 做简单 benchmark. }

\PY{l+s+sd}{    注意:}
\PY{l+s+sd}{    \PYZhy{} 这不是严格的微架构性能实验工具；}
\PY{l+s+sd}{    \PYZhy{} 用于课堂比较“数量级”和相对趋势；}
\PY{l+s+sd}{    \PYZhy{} 正式性能报告应该固定 CPU、电源策略、数据、线程数和环境. }
\PY{l+s+sd}{    \PYZdq{}\PYZdq{}\PYZdq{}}
    \PY{n}{records} \PY{o}{=} \PY{p}{[}\PY{p}{]}

    \PY{k}{for} \PY{n}{name}\PY{p}{,} \PY{n}{func} \PY{o+ow}{in} \PY{n}{cases}\PY{o}{.}\PY{n}{items}\PY{p}{(}\PY{p}{)}\PY{p}{:}
        \PY{k}{for} \PY{n}{\PYZus{}} \PY{o+ow}{in} \PY{n+nb}{range}\PY{p}{(}\PY{n}{warmup}\PY{p}{)}\PY{p}{:}
            \PY{n}{func}\PY{p}{(}\PY{p}{)}

        \PY{n}{gc}\PY{o}{.}\PY{n}{collect}\PY{p}{(}\PY{p}{)}

        \PY{n}{times} \PY{o}{=} \PY{n}{np}\PY{o}{.}\PY{n}{asarray}\PY{p}{(}
            \PY{n}{timeit}\PY{o}{.}\PY{n}{repeat}\PY{p}{(}
                \PY{n}{stmt}\PY{o}{=}\PY{n}{func}\PY{p}{,}
                \PY{n}{repeat}\PY{o}{=}\PY{n}{repeat}\PY{p}{,}
                \PY{n}{number}\PY{o}{=}\PY{n}{number}\PY{p}{,}
            \PY{p}{)}\PY{p}{,}
            \PY{n}{dtype}\PY{o}{=}\PY{n}{np}\PY{o}{.}\PY{n}{float64}\PY{p}{,}
        \PY{p}{)} \PY{o}{/} \PY{n}{number}

        \PY{n}{records}\PY{o}{.}\PY{n}{append}\PY{p}{(}
            \PY{p}{\PYZob{}}
                \PY{l+s+s2}{\PYZdq{}}\PY{l+s+s2}{name}\PY{l+s+s2}{\PYZdq{}}\PY{p}{:} \PY{n}{name}\PY{p}{,}
                \PY{l+s+s2}{\PYZdq{}}\PY{l+s+s2}{median\PYZus{}ms}\PY{l+s+s2}{\PYZdq{}}\PY{p}{:} \PY{n}{np}\PY{o}{.}\PY{n}{median}\PY{p}{(}\PY{n}{times}\PY{p}{)} \PY{o}{*} \PY{l+m+mi}{1000}\PY{p}{,}
                \PY{l+s+s2}{\PYZdq{}}\PY{l+s+s2}{min\PYZus{}ms}\PY{l+s+s2}{\PYZdq{}}\PY{p}{:} \PY{n}{times}\PY{o}{.}\PY{n}{min}\PY{p}{(}\PY{p}{)} \PY{o}{*} \PY{l+m+mi}{1000}\PY{p}{,}
                \PY{l+s+s2}{\PYZdq{}}\PY{l+s+s2}{max\PYZus{}ms}\PY{l+s+s2}{\PYZdq{}}\PY{p}{:} \PY{n}{times}\PY{o}{.}\PY{n}{max}\PY{p}{(}\PY{p}{)} \PY{o}{*} \PY{l+m+mi}{1000}\PY{p}{,}
            \PY{p}{\PYZcb{}}
        \PY{p}{)}

    \PY{k}{return} \PY{p}{(}
        \PY{n}{pd}\PY{o}{.}\PY{n}{DataFrame}\PY{p}{(}\PY{n}{records}\PY{p}{)}
        \PY{o}{.}\PY{n}{sort\PYZus{}values}\PY{p}{(}\PY{l+s+s2}{\PYZdq{}}\PY{l+s+s2}{median\PYZus{}ms}\PY{l+s+s2}{\PYZdq{}}\PY{p}{)}
        \PY{o}{.}\PY{n}{reset\PYZus{}index}\PY{p}{(}\PY{n}{drop}\PY{o}{=}\PY{k+kc}{True}\PY{p}{)}
    \PY{p}{)}


\PY{k}{def}\PY{+w}{ }\PY{n+nf}{plot\PYZus{}benchmark}\PY{p}{(}\PY{n}{result}\PY{p}{:} \PY{n}{pd}\PY{o}{.}\PY{n}{DataFrame}\PY{p}{,} \PY{n}{title}\PY{p}{:} \PY{n+nb}{str}\PY{p}{)} \PY{o}{\PYZhy{}}\PY{o}{\PYZgt{}} \PY{k+kc}{None}\PY{p}{:}
    \PY{n}{ordered} \PY{o}{=} \PY{n}{result}\PY{o}{.}\PY{n}{sort\PYZus{}values}\PY{p}{(}\PY{l+s+s2}{\PYZdq{}}\PY{l+s+s2}{median\PYZus{}ms}\PY{l+s+s2}{\PYZdq{}}\PY{p}{,} \PY{n}{ascending}\PY{o}{=}\PY{k+kc}{True}\PY{p}{)}

    \PY{n}{plt}\PY{o}{.}\PY{n}{figure}\PY{p}{(}\PY{n}{figsize}\PY{o}{=}\PY{p}{(}\PY{l+m+mi}{9}\PY{p}{,} \PY{l+m+mf}{4.5}\PY{p}{)}\PY{p}{)}
    \PY{n}{plt}\PY{o}{.}\PY{n}{barh}\PY{p}{(}\PY{n}{ordered}\PY{p}{[}\PY{l+s+s2}{\PYZdq{}}\PY{l+s+s2}{name}\PY{l+s+s2}{\PYZdq{}}\PY{p}{]}\PY{p}{,} \PY{n}{ordered}\PY{p}{[}\PY{l+s+s2}{\PYZdq{}}\PY{l+s+s2}{median\PYZus{}ms}\PY{l+s+s2}{\PYZdq{}}\PY{p}{]}\PY{p}{)}
    \PY{n}{plt}\PY{o}{.}\PY{n}{xlabel}\PY{p}{(}\PY{l+s+s2}{\PYZdq{}}\PY{l+s+s2}{Median time (ms)}\PY{l+s+s2}{\PYZdq{}}\PY{p}{)}
    \PY{n}{plt}\PY{o}{.}\PY{n}{ylabel}\PY{p}{(}\PY{l+s+s2}{\PYZdq{}}\PY{l+s+s2}{\PYZdq{}}\PY{p}{)}
    \PY{n}{plt}\PY{o}{.}\PY{n}{title}\PY{p}{(}\PY{n}{title}\PY{p}{)}
    \PY{n}{plt}\PY{o}{.}\PY{n}{tight\PYZus{}layout}\PY{p}{(}\PY{p}{)}
    \PY{n}{plt}\PY{o}{.}\PY{n}{show}\PY{p}{(}\PY{p}{)}


\PY{k}{def}\PY{+w}{ }\PY{n+nf}{dataframe\PYZus{}memory\PYZus{}mb}\PY{p}{(}\PY{n}{df}\PY{p}{:} \PY{n}{pd}\PY{o}{.}\PY{n}{DataFrame}\PY{p}{)} \PY{o}{\PYZhy{}}\PY{o}{\PYZgt{}} \PY{n+nb}{float}\PY{p}{:}
    \PY{k}{return} \PY{n}{df}\PY{o}{.}\PY{n}{memory\PYZus{}usage}\PY{p}{(}\PY{n}{index}\PY{o}{=}\PY{k+kc}{True}\PY{p}{,} \PY{n}{deep}\PY{o}{=}\PY{k+kc}{True}\PY{p}{)}\PY{o}{.}\PY{n}{sum}\PY{p}{(}\PY{p}{)} \PY{o}{/} \PY{l+m+mi}{1024}\PY{o}{*}\PY{o}{*}\PY{l+m+mi}{2}


\PY{k}{def}\PY{+w}{ }\PY{n+nf}{file\PYZus{}size\PYZus{}mb}\PY{p}{(}\PY{n}{path}\PY{p}{:} \PY{n}{Path}\PY{p}{)} \PY{o}{\PYZhy{}}\PY{o}{\PYZgt{}} \PY{n+nb}{float}\PY{p}{:}
    \PY{k}{return} \PY{n}{path}\PY{o}{.}\PY{n}{stat}\PY{p}{(}\PY{p}{)}\PY{o}{.}\PY{n}{st\PYZus{}size} \PY{o}{/} \PY{l+m+mi}{1024}\PY{o}{*}\PY{o}{*}\PY{l+m+mi}{2}
\end{Verbatim}
\end{tcolorbox}

    \subsection{5. NumPy:真正理解
ndarray}\label{numpyux771fux6b63ux7406ux89e3-ndarray}

通过这部分希望大家可以理解下面的问题:

\begin{Shaded}
\begin{Highlighting}[]
\NormalTok{ndarray ≠ 更快的 list}

\NormalTok{ndarray =}
\NormalTok{    data buffer}
\NormalTok{    + dtype}
\NormalTok{    + shape}
\NormalTok{    + strides}
\end{Highlighting}
\end{Shaded}

能够解释:

\begin{itemize}
\tightlist
\item
  \texttt{dtype}
\item
  \texttt{shape}
\item
  \texttt{ndim}
\item
  \texttt{size}
\item
  \texttt{itemsize}
\item
  \texttt{nbytes}
\item
  \texttt{strides}
\item
  \texttt{flags}
\item
  \texttt{view}
\item
  \texttt{copy}
\item
  \texttt{reshape}
\item
  \texttt{transpose}
\item
  \texttt{axis}
\end{itemize}

\subsubsection{ndarray 是什么}\label{ndarray-ux662fux4ec0ux4e48}

ndarray 是 Numpy 中的一个基础类型, 它底层的简化定义如下:

\begin{Shaded}
\begin{Highlighting}[]
\KeywordTok{typedef} \KeywordTok{struct} \OperatorTok{\{}
\NormalTok{    PyObject\_HEAD}

    \CommentTok{/* Number of dimensions */}
    \DataTypeTok{int}\NormalTok{ nd}\OperatorTok{;}
    \CommentTok{/* Shape of the array */}
\NormalTok{    npy\_intp }\OperatorTok{*}\NormalTok{dimensions}\OperatorTok{;}
    \CommentTok{/* Strides of the array */}
\NormalTok{    npy\_intp }\OperatorTok{*}\NormalTok{strides}\OperatorTok{;}
    \CommentTok{/* Pointer to the actual data */}
    \DataTypeTok{void} \OperatorTok{*}\NormalTok{data}\OperatorTok{;}
    \CommentTok{/* Data type descriptor */}
\NormalTok{    PyArray\_Descr }\OperatorTok{*}\NormalTok{descr}\OperatorTok{;}
    \CommentTok{/* Base object, used for views */}
\NormalTok{    PyObject }\OperatorTok{*}\NormalTok{base}\OperatorTok{;}
    \CommentTok{/* Flags */}
    \DataTypeTok{int}\NormalTok{ flags}\OperatorTok{;}
    \CommentTok{/* Weak references */}
\NormalTok{    PyObject }\OperatorTok{*}\NormalTok{weakreflist}\OperatorTok{;}
    \CommentTok{/* ... other internal fields ... */}

\OperatorTok{\}}\NormalTok{ PyArrayObject}\OperatorTok{;}
\end{Highlighting}
\end{Shaded}

NumPy 官方 \texttt{ndarray} 文档把
\texttt{shape}、\texttt{strides}、\texttt{ndim}、\texttt{data}、\texttt{itemsize}、\texttt{nbytes}、\texttt{base}、\texttt{dtype}
等直接列为数组(ndarray)的核心属性；\texttt{strides}
表示沿每个维度前进一步需要跨过的字节数.

\begin{Shaded}
\begin{Highlighting}[]
\NormalTok{flowchart LR}
\NormalTok{    A["ndarray"] {-}{-}\textgreater{} B["data buffer"]}
\NormalTok{    A {-}{-}\textgreater{} C["dtype"]}
\NormalTok{    A {-}{-}\textgreater{} D["shape"]}
\NormalTok{    A {-}{-}\textgreater{} E["strides"]}

\NormalTok{    B {-}{-}\textgreater{} B1["真正的数据"]}
\NormalTok{    C {-}{-}\textgreater{} C1["一个元素怎样解释"]}
\NormalTok{    D {-}{-}\textgreater{} D1["逻辑维度"]}
\NormalTok{    E {-}{-}\textgreater{} E1["逻辑坐标如何映射到内存"]}
\end{Highlighting}
\end{Shaded}

所以 \texttt{ndarray} 本质上可以理解为:

\begin{Shaded}
\begin{Highlighting}[]
\NormalTok{Python list}

\NormalTok{┌─────┐}
\NormalTok{│ ptr │──→ Python object}
\NormalTok{├─────┤}
\NormalTok{│ ptr │──→ Python object}
\NormalTok{├─────┤}
\NormalTok{│ ptr │──→ Python object}
\NormalTok{└─────┘}

\NormalTok{NumPy ndarray}

\NormalTok{metadata:}
\NormalTok{shape   = (4,)}
\NormalTok{dtype   = float64}
\NormalTok{strides = (8,)}

\NormalTok{data buffer:}
\NormalTok{┌────────┬────────┬────────┬────────┐}
\NormalTok{│ 1.2    │ 3.4    │ 5.6    │ 7.8    │}
\NormalTok{└────────┴────────┴────────┴────────┘}
\end{Highlighting}
\end{Shaded}

\paragraph{新建一个 ndarray}\label{ux65b0ux5efaux4e00ux4e2a-ndarray}

我们创建一个最基本的 ndarray 如下:

    \begin{tcolorbox}[breakable, size=fbox, boxrule=1pt, pad at break*=1mm,colback=cellbackground, colframe=cellborder]
\prompt{In}{incolor}{ }{\boxspacing}
\begin{Verbatim}[commandchars=\\\{\}]

\end{Verbatim}
\end{tcolorbox}


    % Add a bibliography block to the postdoc
    
    
    
\end{document}
